% Header, centre. \cgTitle is the title submitted to CARTOGRASS 2026, word
% for word (content/cg/text_header.tex), set at 64 pt in two lines broken
% after "intertidal" (\cgTitleBreak; 409 and 491 mm in the 507 mm measure),
% so it leads the 44 pt section headings by a clear step. The fallback
% repeats it.
% Round 3 (2026-10-19): \parskip 0. The panel's default \parskip (0.38
% \baselineskip, about 4.5 mm) was added at each of the three paragraph breaks
% on top of the explicit gaps, so the block took 107.6 of the 118-mm band and,
% centred, its title stood 5 mm under the top edge of the paper (the logo
% frames start 14-16 mm down, the footer keeps 12 mm clear). The gaps are now
% the only space between the lines: the block is at most 96 mm, so centred in
% the band it starts at least 11 mm below the edge.
\setlength{\parskip}{0pt}%
{\color{white}\fontsize{64}{72}\selectfont\bfseries
 \def\cgTitleBreak{\\}%
 \cguse{\cgTitle}{MAREA: a method for mapping intertidal\\topography using
 Sentinel-2 openEO data cubes}\par}
\vspace{5mm}
{\color{PMist}\fontsize{33}{40}\selectfont\itshape
 \cguse{\cgKicker}{The tide arrives late: MAREA reads the delay from
 satellite images and corrects the maps for it.}\par}
\vspace{6mm}
% TODO(authors): names, order and affiliations to be confirmed by Jorge (see
% the note in content/cg/text_header.tex).
{\color{white}\fontsize{30}{36}\selectfont\bfseries
 \cguse{\cgAuthors}{Jorge Al\'{\i}as \'Alvarez\textsuperscript{1},
 Fernando Garc\'{\i}a-Gonz\'alez\textsuperscript{2},
 Pablo Gonz\'alez\textsuperscript{1}}\par}
\vspace{2mm}
{\color{PMist}\fontsize{23}{28}\selectfont
 \cguse{\cgAffilA}{\textsuperscript{1}University of Oviedo, Spain}\qquad
 \cguse{\cgAffilB}{\textsuperscript{2}Instituto Espa\~nol de
 Oceanograf\'{\i}a (IEO-CSIC), Spain}\qquad
 \cguse{\cgContact}{jorge.alias.alvarez@gmail.com}\par}
